%% =============================================================================
%% Pollity.in — What We Are Building (One-Pager)
%% =============================================================================
\documentclass[11pt, a4paper]{article}

\usepackage[T1]{fontenc}
\usepackage[utf8]{inputenc}
\usepackage{lmodern}
\usepackage[margin=2.5cm, top=2cm, bottom=2cm]{geometry}
\usepackage{xcolor}
\usepackage{titlesec}
\usepackage[most]{tcolorbox}
\usepackage{enumitem}
\usepackage{parskip}
\usepackage{microtype}
\usepackage{tabularx}
\usepackage{booktabs}
\usepackage{hyperref}
\usepackage{fancyhdr}
\usepackage{amssymb}

%% ── Colours ──────────────────────────────────────────────────────────────────
\definecolor{pollityblue}{RGB}{31, 97, 141}
\definecolor{pollitygold}{RGB}{212, 172, 13}
\definecolor{lightgray}{RGB}{245, 245, 245}
\definecolor{darktext}{RGB}{30, 30, 30}

%% ── Section Formatting ───────────────────────────────────────────────────────
\titleformat{\section}{\large\bfseries\color{pollityblue}}{}{0em}{}[\color{pollityblue}\titlerule]
\titleformat{\subsection}{\normalsize\bfseries\color{pollityblue}}{}{0em}{}

%% ── Header / Footer ──────────────────────────────────────────────────────────
\pagestyle{fancy}
\fancyhf{}
\lhead{\color{pollityblue}\textbf{Pollity.in}}
\rhead{\color{gray}\small What We Are Building}
\cfoot{\color{gray}\small pollity.in \textbar{} March 2026}
\renewcommand{\headrulewidth}{0.4pt}

\hypersetup{colorlinks=true, urlcolor=pollityblue, linkcolor=pollityblue}

%% =============================================================================
\begin{document}

%% ── Header ───────────────────────────────────────────────────────────────────
\begin{center}
  {\color{pollityblue}\rule{\linewidth}{1.5pt}}\\[0.5em]
  {\Huge\bfseries\color{pollityblue} Pollity.in}\\[0.3em]
  {\large A Digital Chief of Staff for India's Elected Representatives}\\[0.3em]
  {\color{pollityblue}\rule{\linewidth}{1.5pt}}
\end{center}

\vspace{0.8em}

%% ── The Problem ──────────────────────────────────────────────────────────────
\section{The Problem}

India has \textbf{3.3 million elected representatives}: MPs, MLAs, municipal councillors, panchayat sarpanchs. Every single one of them is overwhelmed.

Each week, hundreds of constituents reach out with problems: a broken road, a blocked pension, a corrupt official, a child out of school. These grievances arrive through WhatsApp messages, phone calls, in-person visits, and handwritten letters. There is no system to manage any of it.

The result:

\begin{itemize}[leftmargin=1.5em, itemsep=3pt]
  \item Grievances get lost between intake and resolution
  \item Citizens hear nothing after they file a complaint
  \item Representatives cannot answer ``how many problems did we solve this month?''
  \item Staff spend their days forwarding messages and chasing departments, not solving problems
\end{itemize}

The gap between ``resolved on paper'' and ``actually resolved'' is enormous. No one is measuring it.

%% ── What We Are Building ─────────────────────────────────────────────────────
\section{What We Are Building}

\begin{tcolorbox}[colback=pollityblue!8, colframe=pollityblue, sharp corners, boxrule=0.6pt]
\textbf{Pollity.in is an AI-powered management platform for elected representatives.}
Think of it as a Digital Chief of Staff that handles the operational complexity of running a constituency office, so the representative can focus on actual governance.
\end{tcolorbox}

\vspace{0.8em}

The first module, \textbf{Jan-Sunwai}, tackles grievance management:

\begin{enumerate}[leftmargin=1.5em, itemsep=4pt]
  \item A constituent sends a WhatsApp message describing their problem
  \item Our AI instantly classifies it by type (infrastructure, healthcare, corruption, etc.) and urgency
  \item The constituent gets an acknowledgement with a unique Reference ID, within seconds
  \item The office staff see every grievance on a single dashboard, assigned, tracked, and updated through to resolution
  \item The representative gets data: how many grievances came in, how many were resolved, what the biggest issues are
\end{enumerate}

\vspace{0.5em}

\textbf{No more lost WhatsApp messages. No more citizens calling back with no answer. No more guessing.}

%% ── The Bigger Vision ────────────────────────────────────────────────────────
\section{The Bigger Vision}

Jan-Sunwai is module one. The full Pollity platform will include:

\begin{tabularx}{\linewidth}{lX}
\toprule
\textbf{Module} & \textbf{What It Does} \\
\midrule
Jan-Sunwai      & Grievance management and constituent service \\
Vidhayak-Mitra  & Legislative research: bill summaries, policy analysis, and talking points \\
Vikas-Tracker   & Track constituency development funds, project status, budget utilisation \\
Nagrik-Setu     & Constituent engagement via social media and mass communication \\
\bottomrule
\end{tabularx}

\vspace{0.5em}
Together, these modules give every elected representative the operational infrastructure that only well-funded national politicians have ever had access to, at a fraction of the cost.

%% ── Why It Matters ───────────────────────────────────────────────────────────
\section{Why It Matters}

India's democracy runs on 3.3 million local representatives. Most of them are first-generation politicians with no staff, no systems, and no data. Pollity gives them infrastructure.

When a sarpanch in rural Maharashtra can show her constituency that she resolved 84\% of grievances filed last quarter, that is accountability. That is what builds trust between citizens and their government at the grassroots level.

%% ── Who We Are ───────────────────────────────────────────────────────────────
\section{Who We Are}

\begin{tabularx}{\linewidth}{lX}
\toprule
\textbf{Manav Mutneja} & Co-Founder \& CEO. Antitrust lawyer and digital governance expert. Leads strategy, partnerships, and stakeholder engagement with government institutions. \\[4pt]
\textbf{Piyush Zaware} & Co-Founder \& CTO. Computational engineer and econometrician, MPP candidate at University of Chicago Harris. Builds the technical infrastructure and data systems. \\[4pt]
\textbf{Joe DiTommaso} & Co-Founder \& COO. Technology and management information systems background. Leads product operations, user acquisition, and go-to-market execution. \\
\bottomrule
\end{tabularx}

\vspace{0.5em}
We are currently building toward a live prototype demo in Spring 2026.

%% ── Where We Are ─────────────────────────────────────────────────────────────
\section{Where We Are Right Now}

\begin{tcolorbox}[colback=lightgray, colframe=pollityblue, sharp corners, boxrule=0.5pt]
\begin{itemize}[leftmargin=1.5em, itemsep=3pt, topsep=2pt]
  \item[$\checkmark$] Core product concept validated through customer discovery interviews
  \item[$\checkmark$] Jan-Sunwai backend built and deployed (FastAPI on Railway)
  \item[$\checkmark$] AI classification pipeline live (Claude Haiku, 8 categories, 4 urgency levels)
  \item[$\checkmark$] Grievance dashboard deployed (Streamlit)
  \item[$\rightarrow$] WhatsApp Business API integration in progress
  \item[$\rightarrow$] First beta users targeted post-SNVC (50 offices, Q3 2026)
\end{itemize}
\end{tcolorbox}

\vspace{1em}
\begin{center}
{\color{pollityblue}\rule{0.5\linewidth}{0.4pt}}\\[0.5em]
{\small Interested? Reach out at \textbf{piyush.zaware@pollity.in} \textbar{} \textbf{manavmutneja@pollity.in} \textbar{} \textbf{joe.ditommaso@pollity.in}}
\end{center}

\end{document}
